%
% 2017hiwslides
%
%  Created by Carter Tazio Schonwald on 2017-08-16.
%   Copyright (c) 2017 . All rights reserved.
%
\documentclass[11pt,reqno]{beamer}

% \usepackage{hyperref}
% \usepackage{url}

\usepackage{enumerate}
\usepackage{microtype}

% \usepackage[parfill]{parskip} % Activate to begin paragraphs with an empty line rather than an indent

% %%% PACKAGES
% \usepackage{booktabs} % for much better looking tables
% \usepackage{array} % for better arrays (eg matrices) in maths
\usepackage{paralist} % very flexible & customisable lists (eg. enumerate/itemize, etc.)
% \usepackage{verbatim} % adds environment for commenting out blocks of text & for better verbatim
% \usepackage{subfig} % make it possible to include more than one captioned figure/table in a single float
% % These packages are all incorporated in the memoir class to one degree or another...


% PAGE DIMENSIONS

% \usepackage[letterpaper,margin=1.5in]{geometry} % to change the page dimensions

% \geometry{letterpaper} % or letterpaper (US) or a5paper or....
% % \geometry{margins=2in} % for example, change the margins to 2 inches all round
% % \geometry{landscape} % set up the page for landscape
% %   read geometry.pdf for detailed page layout information
%layout
% \usepackage{fullpage}

%math symbols
% \usepackage{amsthm}
\usepackage{amsmath}
\usepackage{amssymb}

%%%font packages

\usepackage[utf8]{inputenc}
% \usepackage[T1]{fontenc}


% \usepackage{lmodern}
%\usepackage{mathptmx}
%\usepackage[charter]{mathdesign}

% \usepackage{fancyhdr} % This should be set AFTER setting up the page geometry
% \pagestyle{fancy} % options: empty , plain , fancy
% be sure to give short title version

% \renewcommand{\headrulewidth}{0pt} % customise the layout...
% \lhead{}\chead{}\rhead{}
% \lfoot{}\cfoot{\thepage}\rfoot{}


\usepackage{txfonts}
\usepackage[pdftex]{graphicx}
%\usepackage[usenames,dvipsnames]{color}
%\usepackage{tikz}
% \usetikzlibrary{arrows,decorations.pathmorphing,backgrounds,positioning,fit,matrix,calc}

% \usepackage{fixltx2e}

% \usepackage[vlined,ruled,linesnumbered]{algorithm2e}
% \usepackage{xfrac}

% \usepackage[normallists,normalsections,normaltitle]{savetrees}


% thm style defaults
%
% \theoremstyle{plain}
% \newtheorem{theorem}{Theorem}[section]
% \newtheorem{corollary}[theorem]{Corollary}
% \newtheorem{lemma}[theorem]{Lemma}
%
% \theoremstyle{definition}
% \newtheorem{conjecture}{Conjecture}[section]
% \newtheorem{definition}{Definition}[section]
% \newtheorem{example}{Example}[section]
%
% \theoremstyle{remark}
% \newtheorem{observation}{Observation}
% \newtheorem{remark}{Remark}

% time stamps to keep me honest
\usepackage{datetime}
\renewcommand{\dateseparator}{-}
\settimeformat{ampmtime}
 \newdateformat{fullisodate}{
\THEYEAR-\twodigit{\THEMONTH}-\twodigit{\THEDAY},
 \shortdayofweekname{\THEDAY}{\THEMONTH}{\THEYEAR}\ at  \currenttime}
\newcommand{\now}{\fullisodate
  \today  \usdate}

\usepackage{minted}

%labelling aids
\newcommand{\PropositionName}[1]{\label{prop:#1}}
\newcommand{\LemmaName}[1]{\label{lma:#1}}
\newcommand{\TheoremName}[1]{\label{thm:#1}}
\newcommand{\FigureName}[1]{\label{fig:#1}}

\newcommand{\PropositionRef}[1]{Proposition~\ref{prop:#1}}
\newcommand{\LemmaRef}[1]{Lemma~\ref{lma:#1}}
\newcommand{\TheoremRef}[1]{Theorem~\ref{thm:#1}}
\newcommand{\FigureRef}[1]{Figure~\ref{fig:#1}}

\newcommand{\abs}[1]{\ensuremath{\vert #1 \vert }}
\newcommand{\Abs}[1]{\ensuremath{\left \vert #1 \right \vert}}

\newcommand{\R}{\ensuremath{\mathbb R}}
\newcommand{\C}{\ensuremath{\mathbb C}}
\newcommand{\Z}{\ensuremath{\mathbb Z}}
\newcommand{\Q}{\ensuremath{\mathbb Q}}
\newcommand{\N}{\ensuremath{\mathbb N}}
\newcommand{\F}{\ensuremath{\mathbb F}}
\newcommand{\Zn}[1]{\ensuremath{\mathbb Z / #1 \mathbb Z}}


\newcommand{\Rea}{\R}
\newcommand{\Int}{\Z}
\newcommand{\Rat}{\Q}
\newcommand{\Cmp}{\C}
\newcommand{\Nat}{\N}


\newcommand{\expec}{\ensuremath{\mathbb E}}
\newcommand{\Expec}[1]{\ensuremath{\expec\Paren{#1}}}
\newcommand{\Var}{\text{Var}\,}
\newcommand{\prob}{\ensuremath{\mathbb P}}
\newcommand{\Prob}[1]{\ensuremath{\prob\Paren{#1}}}
\newcommand{\M}{\mathcal }
\newcommand{\set}[1]{\ensuremath{\lbrace #1 \rbrace}}
\newcommand{\Set}[1]{\ensuremath{\left \lbrace #1 \right\rbrace}}

\newcommand{\mcal}[1]{\ensuremath{\mathcal{#1}}}

\newcommand{\Paren}[1]{\ensuremath{\left ( #1 \right )}}
\newcommand{\paren}[1]{\ensuremath{( #1 )}}

\newcommand{\floor}[1]{\ensuremath{\lfloor #1 \rfloor}}
\newcommand{\Floor}[1]{\ensuremath{\left \lfloor #1 \right\rfloor}}

\newcommand{\Ceil}[1]{\ensuremath{\left \lceil #1 \right \rceil }}
\newcommand{\ceil}[1]{\ensuremath{\lceil #1 \rceil}}

\newcommand{\poly}{\text{poly}}

\newcommand{\eps}{\ensuremath{\epsilon}}

\newcommand{\inv}[1]{\ensuremath{#1^{-1}}}
\newcommand{\invp}[1]{\ensuremath{\Paren{#1}^{-1}}}

\newcommand{\norm}[1]{\ensuremath{\lVert #1\rVert}}
\newcommand{\Norm}[1]{\ensuremath{\left\lVert #1\right\rVert}}

\newcommand{\inner}[2]{\ensuremath{\left\langle #1,#2 \right\rangle}}
\newcommand{\kernel}[1]{\ensuremath{\text{ker}\Paren{#1}}}
\newcommand{\range}[1]{\ensuremath{\text{ran}\Paren{#1}}}

\newcommand{\card}[1]{\ensuremath{\abs{#1}}}
%\newcommand{\set}[1]{\left \{ #1 \right \}}                     % Set notation: { ... }
\newcommand{\setst}[2]{\ensuremath{\left \{ #1 \mid #2 \right \}}}           % Set notation: { ... | ... }
\newcommand{\union}{\ensuremath{\cup}}
\newcommand{\symdiff}{\ensuremath{\bigtriangleup}}

\newcommand{\bigoh}[1]{\ensuremath{\mcal O \Paren{#1}}}
\newcommand{\bigtheta}[1]{\ensuremath{\varTheta \Paren{#1}}}
\newcommand{\bigomega}[1]{\ensuremath{\Omega \Paren{#1}}}

\newcommand{\liloh}[1]{\ensuremath{o\Paren{#1}}}
\newcommand{\lilomega}[1]{\ensuremath{\omega\Paren{#1}}}

\begin{document}
\title{GHC Core and Linear Logic should be best friends \\ Haskell Implementors}
\author{Carter Tazio Schonwald \& Joel Burget}
\date{created: 2017-08-16\\ presented 2017-09-19 \\ last generated \now}

\maketitle

\begin{frame}
  \frametitle{Introduction - Authors}

\begin{centering}
\begin{tabular}{l|r}
Joel & Carter  \\
  \includegraphics[scale=0.06]{../joelpics/joel.jpg}   &
  \includegraphics[scale=0.22]{../joelpics/IMG_4133.png}
\end{tabular}
\end{centering}

Our opinions are our own. etc etc :)
\end{frame}

\begin{frame}\frametitle{What is context of our work?}
\begin{enumerate}[(i)]
  \item Our work is motivated by the need to design database and software systems
  that support modeling ownership and authorization correctly.
  \item This could be called ``Blockchain for finance/enterprise\ldots done right''
  \item And our particular approach is around developing a functional model
  of linear logic, applying the compiler engineering knowledge ghc and related
  systems embody, and our work is about building a programming language
  that prevents errors in ownership descriptions.
\end{enumerate}

\end{frame}

\begin{frame}\frametitle{What we're here to share today}
  \begin{enumerate}[(i)]
    \item Motivate linearity/ownership with some simple but nontrivial examples of ownership modeled
    in Haskell. We only need the idea of a linear pair for this.
    \item Show how catchable pure exceptions break type safety/soundness of linearity.
    \item Explain the 6 different type formers of classical linear logic in terms of
    their functional haskell analogues.
    \item Show how some of these type formers tie into current work and design directions
    of ghc!
  \end{enumerate}
\end{frame}


% \begin{frame}[t]\frametitle{title}



% \end{frame}

% what is linear logic (high level)

% heres an example (cake or gum)

% linear logic for  haskellers

% how linear logic ideas intersect with design activity in ghc / ghc core




\begin{frame}[fragile]
\frametitle{What \emph{is} a Loan?}

\begin{minted}{haskell}
import Control.Concurrent.STM
import Control.Concurrent.STM.TQueue

newtype InQ  a = InQ  { inQ  :: TQueue a }
newtype OutQ a = OutQ { outQ :: TQueue a }

type Repay = InQ BubbleGum
type Debt = OutQ BubbleGum

data Loan = Loan
  { principal :: Chocolate
  , repay     :: Repay }

-- So many offers. Do I want to take out this loan?
assessLoan :: Loan -> TVar (Set Loan) -> STM (Maybe Loan)
assessLoan loan wallet = undefined
\end{minted}
\end{frame}

% measured monoids sound nice for aggregating payments
% integer vs fractional fungibible

\begin{frame}
  \frametitle{What can I do with my bubble gum OutQueue's?}
Alice and Bob both owe me some bubblegum. Rick (who's very risk averse) will happily pay me
for the first 10 pieces of gum I receive, and his friend Morty (always taking risks) will
also pay me for any gum after the first 10! In some circles, this sort of rearrangement
is called a \emph{Sequential Tranche}. Sounds like a great job for Concurrency!


So lets figure out how to recompose the OutQueues from Alice and Bob into new
OutQueues which Rick and Morty can pay me for.

\end{frame}

\begin{frame}[fragile]
\frametitle{Concurrency Composes Ownership (1/2)}

\begin{minted}{haskell}
newDebt :: IO (Repay, Debt)
newDebt = fmap (InQ &&& OutQ) newTQueueIO

tranche :: Debt -> Debt -> IO (Debt, Debt)
tranche alice bob = do
    (mortyIn, morty) <- newDebt
    (rickIn, rick) <- newDebt
    _ <- forkIO (allocate alice bob mortyIn rickIn 10)
    return (morty, rick)

allocate :: Debt -> Debt -> Repay -> Repay -> Int -> IO ()
allocate alice bob morty rick remain = do
    taken <- atomically (divvy alice bob morty rick remain)
    allocate alice bob morty rick (remain - taken)
\end{minted}
\end{frame}

\begin{frame}[fragile]
\frametitle{Concurrency Composes Ownership (2/2)}

\begin{minted}{haskell}
divvy :: Debt -> Debt -> Repay -> Repay -> Int -> STM Int
divvy alice bob morty rick remain = do
    ma <- tryReadTQueue (outQ alice)
    mb <- tryReadTQueue (outQ bob)
    let available = catMaybes [ma, mb]
    let (safe, risky) = splitAt remain available
    mapM_ (writeTQueue (inQ morty)) safe
    mapM_ (writeTQueue (inQ rick)) risky
    return (length available)
\end{minted}
\end{frame}

% \begin{frame}
%   \frametitle{Motivation}
%   Example of bad core
% \end{frame}

\begin{frame}[fragile]
\frametitle{linear typing is unsound in the presence of catchable exceptions implicit }
% Consider this innocent program, \ldots are we cartesian closed or merely (symmetric) braided monoidal?
% this should be 2-3 slides in after 1-2 terms/ideas some intro stuff
% \begin{onlyenv}<1>
\begin{minted}{haskell}
module BadUncurry where
import Control.Exception

data Fst a = Only a deriving (Show, Typeable, Exception)

badFst :: (Typeable a, Show a) => a -> b -> c
badFst a = throw (Only a)

notOnly (Only a) = a

nonLinear :: (Show a, Typeable a) => (a, b) -> IO a
nonLinear pr = catch
  (return  (uncurry badFst pr))
  (return . notOnly)
\end{minted}
\end{frame}

\begin{frame}\frametitle{Possible strategies to recover soundness?}
Linear types, even just pairs \ldots seem to be unsound in the
face of implicit effects.

\begin{enumerate}[(i)]
  \item Type and Effect systems!  % solves the implicit effects
  \item restrict how pure exceptions can be caught.
  \item only consider exceptions in an STM / transactional effects with rollback setting.
        (this is the semantics we're delving into for our work).
  \item Perhaps some sort of checkpointing style roll back as seen
\end{enumerate}

These strategies all seem like an ill fit for GHC today.

\end{frame}


\begin{frame}
  \frametitle{Back to business: Linear Logic for the Working Haskeller}
The big  6 type operators in linear logic are:

\begin{tabular}{c|r|c|l}
Linear & Haskell & HS CPS (\neg \neg) & Name \\
\hline
\( \neg X \) & \( \neg X \) \emph{or} \( X \rightarrow \forall a.a  \) &  same   &  negation \\
\( A \otimes B \) & \( (A,B) \) &  & tensor or product \\
\( A \oplus B \) & \( A + B  \) & & sum or either \\
\( A \& B \) & ? & \( \neg (\neg A + \neg B) \) & With or (External) Choice \\
\( A \invamp B \) & ? & \( \neg (\neg A , \neg B) \) & Par \\

\end{tabular}

Aside: If you start with a function type that has unboxed tuple arguments/results
and sum types, you can encode all of these in the CPS style. See eg: types are calling conventions.
% \(\invamp   \)
\end{frame}

\begin{frame}\frametitle{We dont always need linearity to get benefits from linear logic}
By relating various control and abstraction representations to their linear logical
siblings, a number of really lovely and potentially useful design clarifications
arise!


\end{frame}

\begin{frame}\frametitle{What is \( \& \) (Choice)}
Its sort of a ``right hand side of case expressions'' that can act as a generalized
join point representation.

Modeling the RHS of a case as a collection of functions may force needless closure
allocation when in fact that all share the same environment! Choice (\&) can
be used to encapsulate a collection of continuations (join points) that work
in the same scope!
\end{frame}


\begin{frame}[fragile]\frametitle{Choice makes case expressions first class}

at least for join points, we can go first class
\begin{minted}{haskell}
case :: ( a + b , a & b ) -> forall r . r
\end{minted}


\end{frame}

\begin{frame}[fragile]\frametitle{What is Par?}
consider the humble forkIO

\begin{minted}{haskell}
forkIO :: IO () -> IO ThreadId
forkIOCPS :: IO ()->(ThreadId -> IO ())-> forall c . c
forkSimpleCPS ::
  (() -> forall c.c ,() -> forall c.c ) -> forall c . c
\end{minted}
hrmmm ... forkSimpleCPS's type can be written \( \neg (\neg (), \neg ()) \),
which is a special case of \( \neg (\neg a , \neg b) \). Par is a generalization of
forkIO sans ThreadID!

And its resource aware!

    %  a sort of generalized sibling to ForkIO!
    %
\end{frame}




\begin{frame}[fragile]\frametitle{Units, so many units}

% \begin{centering}
\begin{tabular}{lc}
 \( \& \) & empty case \\ \hline
\( \oplus \) & Void type \\ \hline
\( \otimes \) & \( () \) aka \(1 \) \\ \hline
\( \invamp \) & close thread / finalize
\end{tabular}
% \begin{centering}


\end{frame}

\begin{frame}\frametitle{GHC Core already has generalized Linear types}

Demand analysis and the usage info on every binding \emph{IS} a very interesting
fancy linear type system! It just lives in a different place than the rest of the
types... and are only used by the optimizer! It has ideas like \( \geq1 \) and \( \leq 1 \)
to name but the simplest two!

linearity / usage here roughly means ``entered the closure'', which is stronger
than some other notions that seem to be used.
(the range of possible notions isn't cleanly articulated anywhere I've seen)

\end{frame}



%

% \begin{frame}[fragile]
% what next? What if we replaced our definition of uncurry?
% \begin{minted}{haskell}
% linearUncurry :: (a,b) -> ((# a , b #) -> c) -> c
% \end{minted}
% this version of uncurry doesn't allow an exceptional definition of first!
% \end{frame}

% \begin{frame}
%   \frametitle{Demand Analysis}
%   What's the relationship to demand analysis?

%   GHC ALREADY has a very precise sort of linear logic, ... usage annotations
% \end{frame}

% \begin{frame}
%   \frametitle{Demand Analysis}
%   Demand analysis has more possible demands:

%   {\large Strictness demands}
%   \begin{itemize}
%     \item B: hyperstrict
%     \item L: lazy
%     \item S: head strict
%     \item S(s1 ... sn): structured product demand
%     \item C(s): call demand
%   \end{itemize}

%   {\large Absence / usage demands}
%   \begin{itemize}
%     \item A: unused
%     \item U: used on some execution path
%     \item H: head-used (value used but components are ignored)
%     \item U(u1 ... un): structured usage demand
%     \item C(u): call-demand
%   \end{itemize}
% \end{frame}

% \begin{frame}
%   \frametitle{Demand Analysis}
%   We end up with demands like \tt{<S,U><S,U(UA)>m}.
% \end{frame}

% \begin{frame}
%   \frametitle{Demand Analysis}
%   But the details aren't important
% \end{frame}

% \begin{frame}
%   \frametitle{Demand Analysis}
%   Perhaps the main difference is that demand analysis is focused on \emph{inferring} usage without annotations, while a linear core would have checked annotations.
% \end{frame}

\end{document}
